%! Multiplying \& Dividing Rational Expressions — Unit 4, Lesson 4.2 — Group Activity (Answer Key)
\documentclass[10pt]{article}
\usepackage{algebra2-article}
\usepackage{algebra2-key}   % pulls in -boxes; do NOT also load -boxes
\newcommand{\TallMath}[1]{$\displaystyle #1\rule[-1.4em]{0pt}{3.2em}$}

\begin{document}

\pageheader{Unit 4, Lesson 4.2}{Group Activity --- Answer Key}
\namepartnerperiod

\vspace{0.08in}

\begin{headlinebox}{forestbg}
\small\textbf{Factor First, Flip Second, Restrict Always.} Every answer today still has \emph{two}
parts --- the simplest form \textbf{and} the restrictions --- but now the restrictions come from
\textbf{three} places: each original denominator, and the \emph{numerator} of anything you divide by.
With your partner, factor before you touch anything, and be ready to say which source each restriction
came from.
\end{headlinebox}

\vspace{0.06in}

\begin{tcolorbox}[colback=white, colframe=black!40, title=\textbf{Tier R --- Remediate},
    fonttitle=\bfseries, arc=1.5mm, left=3mm, right=3mm, top=2mm, bottom=2mm, breakable]
Simplify each expression and list \emph{all} restrictions. The first one is started for you.

\vspace{4pt}
\begin{enumerate}[leftmargin=1.9em, itemsep=9pt]
  \item $\dfrac{3x}{4}\cdot\dfrac{8}{9x^2}$: \; multiply across to get $\dfrac{24x}{36x^2}$, then
        divide out $12$ and one $x$.

        Simplest form: \ans{$\frac{2}{3x}$} \quad Restriction: $x\neq$ \ans{0}
  \item $\dfrac{x+2}{x-1}\cdot\dfrac{x-1}{x+5}$ \\[3pt]
        Divide out: \ans{$(x-1)$} \quad Simplest form: \ans{$\frac{x+2}{x+5}$} \quad
        $x\neq$ \ans{$1,\ -5$}
  \item $\dfrac{x^2-4}{x+3}\cdot\dfrac{2}{x-2}$ \\[3pt]
        Factored: \ans{$\frac{(x-2)(x+2)}{x+3}\cdot\frac{2}{x-2}$} \quad Simplest form: \ans{$\frac{2(x+2)}{x+3}$} \quad
        $x\neq$ \ans{$-3,\ 2$}
  \item $\dfrac{x}{x+1}\div\dfrac{x-3}{x+1}$ \\[3pt]
        Keep--change--flip: \ans{$\frac{x}{x+1}\cdot\frac{x+1}{x-3}$} \quad Simplest form: \ans{$\frac{x}{x-3}$} \\[3pt]
        Restrictions: from the denominators, $x\neq$ \ans{$-1$}; \; from the divisor's
        \emph{numerator}, $x\neq$ \ans{3}
  \item In problem 4, $x=3$ was never a denominator in the problem as written. Why is it excluded
        anyway? \ans{at $x=3$ the divisor equals $0$, and you cannot divide by $0$}
\end{enumerate}
\end{tcolorbox}

\begin{tcolorbox}[colback=white, colframe=black!40, title=\textbf{Tier A --- Approaching Proficiency},
    fonttitle=\bfseries, arc=1.5mm, left=3mm, right=3mm, top=2mm, bottom=2mm, breakable]
\textbf{Part 1 --- The full table.} Fill in every cell. Two of these are divisions, one hides a pair of
\emph{opposite} binomials, and one has an answer that will surprise you.
\begin{center}
\renewcommand{\arraystretch}{1.7}
\begin{tabularx}{\linewidth}{@{}l X X X@{}}
\hline
\textbf{Expression} & \textbf{Factored form} & \textbf{Restrictions} & \textbf{Simplest form} \\
\hline
\TallMath{\frac{10x^3}{3y^2}\cdot\frac{9y^5}{5x}} & \ans{$\frac{5x\cdot2x^2}{3y^2}\cdot\frac{3y^2\cdot3y^3}{5x}$} & \ans{$x\neq0,\ y\neq0$} & \ans{$6x^2y^3$} \\
\TallMath{\frac{x^2+5x+6}{x^2-9}\cdot\frac{x-3}{x+2}} & \ans{$\frac{(x+2)(x+3)}{(x-3)(x+3)}\cdot\frac{x-3}{x+2}$} & \ans{$x\neq3,\ -3,\ -2$} & \ans{$1$} \\
\TallMath{\frac{x^2-2x-8}{x^2-16}\div\frac{x+2}{x^2+4x}} & \ans{$\frac{(x-4)(x+2)}{(x-4)(x+4)}\cdot\frac{x(x+4)}{x+2}$} & \ans{$x\neq4,\ -4,\ 0,\ -2$} & \ans{$x$} \\
\TallMath{\frac{9-x^2}{x+4}\cdot\frac{x^2+4x}{x-3}} & \ans{$\frac{-(x-3)(x+3)}{x+4}\cdot\frac{x(x+4)}{x-3}$} & \ans{$x\neq-4,\ 3$} & \ans{$-x(x+3)$} \\
\hline
\end{tabularx}
\end{center}

\vspace{4pt}
\textbf{Part 2 --- Error analysis.} A student turns in this work:
\begin{center}
\fbox{\parbox{0.66\linewidth}{\centering \vspace{2pt}
$\dfrac{x}{4}\div\dfrac{5}{x} \;=\; \dfrac{1}{4}\div\dfrac{5}{1} \;=\; \dfrac{1}{20}$ \\[3pt]
{\small ``The two $x$'s were on top and bottom, so I cancelled them first.''}\vspace{2pt}}}
\end{center}
\begin{enumerate}[label=(\alph*), leftmargin=1.9em, itemsep=6pt]
  \item Disprove it with a number. At $x=2$: the original equals \ans{$\frac{1}{5}$} and the student's
        answer equals \ans{$\frac{1}{20}$}.
  \item Do it correctly: $\dfrac{x}{4}\div\dfrac{5}{x}=$ \ans{$\frac{x^2}{20}$}, \; $x\neq$ \ans{0}.
        Check it at $x=2$: \ans{$\frac{4}{20}=\frac15$}
  \item State the rule in one sentence, so the student never does it again.
        \ansline{Flip first, then cancel. The $\div$ symbol is not a fraction bar, so nothing is lined up}
        \ansline{across it --- only after keep--change--flip do numerators and denominators actually face each other.}
\end{enumerate}
\end{tcolorbox}

\begin{tcolorbox}[colback=white, colframe=black!40, title=\textbf{Tier E --- Extension},
    fonttitle=\bfseries, arc=1.5mm, left=3mm, right=3mm, top=2mm, bottom=2mm, breakable]
\textbf{Part 1 --- Restriction detective.} Simplify
\[ \frac{x^2-1}{x^2+5x+6}\div\frac{x^2-x}{x+3} \]
Simplest form: \ans{$\frac{x+1}{x(x+2)}$} \quad Complete restriction list: $x\neq$ \ans{$0,\ 1,\ -2,\ -3$}

\vspace{3pt}
Now say where each one came from. Check every column that applies.
\begin{center}
\renewcommand{\arraystretch}{1.5}
\begin{tabularx}{\linewidth}{@{}c X X X X@{}}
\hline
\textbf{Value} & \textbf{First denominator} & \textbf{Divisor's denominator} & \textbf{Divisor's numerator} & \textbf{Visible in the answer?} \\
\hline
$x=0$  & \ans{---} & \ans{---} & \ans{yes} & \ans{yes} \\
$x=1$  & \ans{---} & \ans{---} & \ans{yes} & \ans{no} \\
$x=-2$ & \ans{yes} & \ans{---} & \ans{---} & \ans{yes} \\
$x=-3$ & \ans{yes} & \ans{yes} & \ans{---} & \ans{no} \\
\hline
\end{tabularx}
\end{center}
One row has \emph{two} sources checked. Which, and why does that not change the answer?
\ansline{$x=-3$ --- it makes both the first denominator and the divisor's denominator zero. A value is either excluded or it is not; being excluded twice is still excluded once.}

\vspace{5pt}
\textbf{Part 2 --- Build one backwards.} Write a \emph{division} problem whose simplified answer is
$\dfrac{x}{x-4}$ and whose restrictions are exactly $x\neq4$, $x\neq-1$, and $x\neq2$.
\begin{center}
\ans{$\frac{x(x-2)}{(x+1)(x-4)}$} $\div$ \ans{$\frac{x-2}{x+1}$}
\end{center}
\begin{enumerate}[label=(\alph*), leftmargin=1.9em, itemsep=5pt]
  \item Which source produced $x\neq2$ in your problem, and how did you make sure that factor still
        divided out of the answer?
  \item Which of your three values is produced by \emph{two} sources at once?
\end{enumerate}
\ansline{(a) The divisor's \emph{numerator}: $(x-2)$ sits on top of the divisor, so it excludes $x=2$; after the flip it lands}
\ansline{in the denominator, where the matching $(x-2)$ upstairs divides it out. (b) $x=-1$ --- it is the first fraction's denominator \emph{and} the divisor's denominator.}

\vspace{5pt}
\textbf{Part 3 --- Recover a length.} A rectangle has area $\dfrac{x^2+7x+12}{x^2}$ square units and
width $\dfrac{x+3}{x}$ units.
\begin{enumerate}[label=(\alph*), leftmargin=1.9em, itemsep=6pt]
  \item Length $=$ area $\div$ width $=$ \ans{$\frac{x+4}{x}$}, \; with restrictions $x\neq$ \ans{$0,\ -3$}.
  \item \textbf{Check by multiplying:} width $\times$ length $=$ \ans{$\frac{(x+3)(x+4)}{x^2}$}. Does it match the
        given area? \ans{Yes}
  \item The restriction $x\neq-3$ came from the divisor's numerator. Substitute $x=-3$ into the
        \emph{width} and explain what it would mean for the rectangle.
        \ansline{the width would be $\frac{0}{-3}=0$ --- a rectangle with no width at all, so there is no length to recover; the algebra's ``you cannot divide by zero'' and the geometry say the same thing}
  \item For a real rectangle both dimensions must be positive. Which $x$-values make sense? Write
        your answer as a union of intervals. \ans{$(-\infty,-4)\cup(0,\infty)$}
\end{enumerate}
\end{tcolorbox}

\begin{teachernote}
\textbf{Tier A, Part 1, row 2} is the row groups argue about: everything divides out and the answer is
literally $\mathbf{1}$ --- yet the expression still carries three restrictions. Push on it: ``if it
equals $1$, what happens at $x=3$?'' (Nothing --- the original has no value there.) Row 3 is today's
showcase: the answer is the single letter $x$, and it carries \emph{four} restrictions, two of them
from the divisor. Row 4 is the Lesson~4.1 sign check; accept $-x(x+3)$, $-(x^2+3x)$, or
$x(-x-3)$, but not a lost $-1$.

\vspace{3pt}
\textbf{Tier A, Part 2} must be done with numbers, not argument. Groups who cancelled across the $\div$
usually believe it because top-with-top \emph{does} happen to work for division; the counterexample is
built on the other diagonal (top of the first against the bottom of the second), which never works.

\vspace{3pt}
\textbf{Tier E, Part 1} is the assessment target of the whole lesson in one table --- if a group can
fill it in, they have the three-source rule. The $x=-3$ row (two sources) and the $x=1$ row (invisible
in the answer, from the divisor's numerator) are the two that separate groups.

\vspace{3pt}
\textbf{Tier E, Part 3(d)} surprises most students: $x<-4$ also gives two positive dimensions, since
both fractions are then negative-over-negative. Accept $(0,\infty)$ with a correct justification, but
push a strong group to test $x=-5$ and find the second interval --- it is a good reminder that
``positive'' has to be checked, not assumed.
\end{teachernote}

\end{document}
